\subsection{可分扩张中的导子}

我们已经看到（V，第 122 页，定理 1），特征为 0 的域 K 的每一个扩张 L 都是可分扩张；此外，K 到 L 上向量空间的每一个导子都能延拓为 L 上的导子（V，第 130 页，推论 1）。更一般地，我们有如下命题：

定理 3. --- 设 K 是一个域，L 是 K 的一个扩域。要使 L 在 K 上可分，其必要且充分条件是 K 到 L-向量空间的每一个导子都能延拓为 L 上的一个导子。

我们可以假设 K 的特征为 \(\neq0\) 。首先假设 L 在 K 上是可分的。设 V 是 L 上的向量空间，且 \(\Delta\) 是从 K 到 V 的一个导子。根据麦克莱恩

判别法（V，第 124 页，注记），域 \(L^p\) 与 K 在 \(K^p\) 上线性不相交。由于 \(\Delta\) 是 K 到 V 的一个 \(K^p\) -线性映射，它以唯一方式延拓为一个 \(L^p\) -线性映射 \(\Delta'\) （\(K[L^p] = K(L^p)\) 到 V 中的）。显然， \(\Delta'\) 是 \(K(L^p)\) 到 V 的一个在 \(L^p\) 上为零的导子；因此它（V，第 101 页，推论 1）延拓为 L 到 V 的一个导子。

反之，假设 K 的每个取值于 L 的导子都能延拓为 L 到 L 的导子。设 B 为 K 的一个 p-基（V，第 99 页），并设 \(\Lambda\) 为由 0 到 p-1 之间的整数组成的、具有有限支撑的族 \((\alpha_{b})_{b\in\mathbf{B}}\) 的集合。对于每个 \(b\in B\) ，存在 K 到 K 的一个导子 \(\Delta_{b}\) ，其特征由 \(\Delta_{b}(b')=\delta_{bb'}\) （克罗内克符号）对每个 \(b'\in B\) 给出（V，第 103 页）。根据假设，对每个 \(b\in B\) ，存在 L 到 L 的一个导子 \(D_{b}\) ，它延拓 \(\Delta_{b}\) 。我们有 \(\mathrm{D}_{b}(b')=\delta_{bb'}\) （对于 b， \(b'\) 属于 B），这证明了在 \(\Omega_{L^{p}}(L)\) 中微分 db（对于 \(b\in B\) ）在 L 上是线性无关的。因此（V，第 102 页，定理 1），B 在 \(L^{p}\) 上是 p-自由的。由此可得（V，第 98 页，命题 1 及第 124 页，注记），K 的扩域 L 是可分的。

推论. --- 若 L 是 K 的可分扩张，则典范映射 \(\alpha_{\mathrm{P}}:\Omega_{\mathrm{P}}(\mathrm{K})\otimes_{\mathrm{K}}\mathrm{L}\to\Omega_{\mathrm{P}}(\mathrm{L})\) 对 K 的每个子域 P 都是单射。反之，若存在 K 的一个完全子域 P（例如 K 的素子域）使得映射 \(\alpha_{P}\) 是单射，则 L 在 K 上是可分的。

第一个断言由命题 2（V，第 128 页）和定理 3 得出。反之，设 P 为 K 的一个完全子域；我们有 \(P = P^{p} \subset K^{p}\) ，因此从 K 到向量 L-空间的每个导子都是 P-导子；推论的第二条断言现在由推论 2（V，第 130 页）和定理 3 得出。

